\subsection{Ambient and history assumptions shift the fitted delay}
Allowing record-specific ambient offsets reduces joint calibration error more
than allowing only an initial history fraction in the representative models
(Fig.~\ref{fig:nuisance-sensitivity}). The fitted delay also changes, although
the cold one-stage formulation retains the lowest calibration error among these
four representatives. With both nuisance terms free, every displayed step-record
ambient offset reaches the imposed upper bound
(Table~\ref{tab:nuisance-parameters}). The lower residual error therefore does
not identify an independently measured ambient temperature. Extending that bound
until the fit looks satisfactory would change the sensitivity question rather
than resolve the missing measurement.
% Generated by report_followups.py; scalar sources in followup_plot_data.json.
\begin{table}[tbp]
\centering\small
\begin{tabular}{lrrrrrr}
\toprule
Model & \shortstack{$K$\\(${}^\circ$C/\%)} & $\theta$ (s) & \shortstack{$\Delta T_{a,\sin}$\\(${}^\circ$C)} & \shortstack{$\Delta T_{a,\rm step}$\\(${}^\circ$C)} & $\gamma_{\sin}$ & $\gamma_{\rm step}$ \\
\midrule
$n=1$ & 0.7390 & 17.159 & 1.339 & 2.000 & 0.000 & 0.004 \\
$n=2$ & 0.7406 & 16.350 & 1.336 & 2.000 & 0.088 & 0.089 \\
$n=12$ & 0.7421 & 15.932 & 1.321 & 2.000 & 0.195 & 0.173 \\
Pure delay & 0.7424 & 15.867 & 1.318 & 2.000 & 0.218 & 0.192 \\
\bottomrule
\end{tabular}
\caption{Joint fits with free ambient and history. Ambient offsets are in degrees Celsius; history fractions $\gamma$ multiply the initial command. Gain units are degrees Celsius per command percentage point. Every displayed step offset reaches the imposed $+2$ degree bound. These are sensitivity fits under declared bounds, not independently measured ambient temperatures or histories.}
\label{tab:nuisance-parameters}
\end{table}


The nuisance fit retains the first sensor observation as its thermal initial
condition while changing the assumed ambient and delay-state history. Its
history fraction describes a declared state family; it does not reconstruct an
unrecorded heater command. The separate constant-input derivation in
Appendix~\ref{sec:nuisance-assumptions} explains why a fitted history or ambient
can also introduce exact parameter ambiguity. A favorable joint residual score
does not establish estimability in every experiment.

\begin{figure}[p]
 \centering
 \includegraphics[width=\linewidth]{nuisance_sensitivity.pdf}
 \caption{Ambient and initial-history assumptions change joint calibration
 error and the selected delay. Cold, fixed is the original linear-hold fit;
 History frees one delay-state history fraction per record with ambient fixed;
 Ambient frees one ambient offset per record with cold delay states; Both frees
 both. Offset bounds are $\pm2\,{}^\circ\mathrm C$ and history fractions lie in
 $[0,1]$. The two thermal experiments share their model parameters. Lines join
 categorical assumptions; markers are fitted representatives without confidence
 intervals. The one-stage delay uses the selected short-delay pole label.}
 \label{fig:nuisance-sensitivity}
\end{figure}

\subsection{Command reconstruction changes the inferred delay}
Changing the command hold shifts the fitted delay while producing smaller
changes in calibration RMSE (Fig.~\ref{fig:reconstruction-sensitivity}). A
backward-Euler fit can even reduce the residual score relative to exact
propagation. That does not make its discretization a more accurate description
of the continuous thermal system: fitted parameters can absorb numerical and
command-reconstruction differences. The recorded time/command pairs do not
identify the command waveform between samples. Consequently, the narrowest
conditional delay intervals should not be interpreted as resolving this
unmeasured waveform choice. All compared holds use the same observations and
the declared cold prehistory; pure delay has no backward-Euler counterpart in
this study.

\begin{figure}[p]
 \centering
 \includegraphics[width=\linewidth]{reconstruction_sensitivity.pdf}
 \caption{Different command reconstructions and propagators change the
 inferred delay at similar residual error. Each point is a separate joint cold
 fit. Linear, left and right denote the assumed hold between recorded command
 values. Exact uses the analytic LTI propagator for that hold; Right Euler uses
 the inherited backward-Euler/right-end convention at the recorded observation
 times. The pure-delay curve ends at the exact right-hold case. Lines join
 categorical assumptions and imply neither a convergence path nor uncertainty
 intervals.}
 \label{fig:reconstruction-sensitivity}
\end{figure}

\subsection{Prediction precision remains conditional on the mean model}
Prediction bands provide a coordinate-invariant target even when the original
physical coordinates have an exact null direction. Its zero-state equivalence
family preserves the measured thermal response, so prediction variance can be
computed on the identifiable quotient without assigning finite uncertainty to
each physical parameter. By contrast, constant-input calibration with
equilibrated delay states leaves a globally free delay that affects a varying
future input. Those target bands are withheld. The saved finite delay-exploration
envelopes are sensitivity plots without a probability interpretation.

For the jointly fitted cold representatives, mean-band widths depend strongly
on the assumed covariance or residual resampling
(Fig.~\ref{fig:prediction-mean-widths}). Fitted IID variance gives narrower mean
bands than the displayed HAC, OU and block-residual alternatives. Wider bands
under another working error law do not establish uncertainty from between-record
mean discrepancy. Table~\ref{tab:prediction-widths} separates mean-band width
from descriptive containment by observation bands on V01. Adding fitted future
error variance does not resolve that record's transfer error. V01 is the first
file in the existing sorted code map, with no selection by residual size; the
complete saved summary retains every validation file, including those with
better containment.
% Generated by report_followups.py; scalar sources in followup_plot_data.json.
\begin{table}[tbp]
\centering\small
\begin{tabular}{lrrr}
\toprule
Model, method & \shortstack{Sine width\\(${}^\circ$C)} & \shortstack{V01 width\\(${}^\circ$C)} & V01 obs. fraction \\
\midrule
$n=1$, IID fitted & 0.078 & 0.061 & 0.135 \\
$n=1$, OU fitted & 0.439 & 0.347 & 0.139 \\
$n=2$, IID fitted & 0.088 & 0.077 & 0.140 \\
$n=2$, OU fitted & 0.544 & 0.479 & 0.154 \\
$n=12$, IID fitted & 0.099 & 0.094 & 0.150 \\
$n=12$, OU fitted & 0.662 & 0.629 & 0.166 \\
Pure delay, IID fitted & 0.102 & 0.097 & 0.161 \\
Pure delay, OU fitted & 0.560 & 0.512 & 0.171 \\
\bottomrule
\end{tabular}
\caption{Median pointwise mean-band widths and descriptive V01 containment by nominal observation-prediction bands. The last column uses observation bands with an added fitted IID or OU marginal future-error variance; it is not containment by the narrower mean bands. V01 is defined in Table~\ref{tab:validation-files}. These fractions on one observed record are not repeated-experiment coverage estimates. HAC and block-residual methods supply no future-error law here and therefore no observation-band containment.}
\label{tab:prediction-widths}
\end{table}


\begin{figure}[p]
 \centering
 \includegraphics[width=\linewidth]{prediction_mean_widths.pdf}
 \caption{Conditional thermal-mean precision depends on the working error
 model. Points give median pointwise nominal 95\% mean-band widths over the
 sine calibration record and held-out V01, excluding their conditioned first
 observations. IID fitted uses residual degrees of freedom on the identifiable
 quotient; HAC60 uses Bartlett bandwidth 60; OU fitted uses the observed
 mean/length Hessian with variance profiled; Block60 uses percentile predictions
 from saved joint circular-block draws. The OU means are separately fitted;
 other methods retain the baseline mean. Lines join methods as visual guides,
 without a common unconditional confidence or repeated-experiment coverage
 interpretation. These are mean bands, not future-observation bands.}
 \label{fig:prediction-mean-widths}
\end{figure}
